\chapter{Probing the driven junction}\label{ch:probes}

A driven junction can be probed by a second, weak perturbation: a probe tone in two-tone
spectroscopy, a tunnelling contact, or a spin-resonance coil. This chapter describes the
transitions that such a probe induces between Floquet states, and how \pkg{} computes
their positions and strengths.

\section{The Floquet golden rule}

Add to the driven Hamiltonian a weak probe $\lambda X(t)\cos\omega_pt$, where $X(t)$ is an
operator, possibly modulated by the drive, with harmonics $X(t)=\sum_kX_k\ee^{\ii k\Omega t}$.
Between Floquet states the probe matrix element oscillates at a ladder of frequencies,
\begin{equation}\label{eq:floquet-element}
  \matel{\psi_\alpha(t)}{X(t)}{\psi_\beta(t)}=\sum_qX^{(q)}_{\alpha\beta}\,
  \ee^{\ii(\qe_\alpha-\qe_\beta+q\Omega)t},\qquad
  X^{(q)}_{\alpha\beta}=\sum_{m,k}\matel{u_{\alpha,m+k-q}}{X_k}{u_{\beta,m}} .
\end{equation}
First-order perturbation theory, averaged over many drive periods, gives the transition rate
from $\beta$ to $\alpha$,
\begin{equation}\label{eq:floquet-golden}
  \Gamma_{\beta\to\alpha}(\omega_p)=\frac{\pi\lambda^2}2\sum_q\abs{X^{(q)}_{\alpha\beta}}^2
  \Bigl[\delta\bigl(\omega_p-\omega^{(q)}_{\alpha\beta}\bigr)+\delta\bigl(\omega_p+\omega^{(q)}_{\alpha\beta}\bigr)\Bigr],
  \qquad \omega^{(q)}_{\alpha\beta}=\qe_\alpha-\qe_\beta+q\Omega ,
\end{equation}
the Floquet generalization of Fermi's golden rule~\cite{grifoni1998driven,kohler2005driven}.
Each pair of Floquet states contributes a ladder of lines separated by the drive frequency.
In the photon picture of \cref{sec:photons}, $q$ counts the drive photons exchanged during
the transition, and $\abs{X^{(q)}_{\alpha\beta}}^2$ is the strength of the $q$-photon line. A
line at a positive frequency absorbs a probe quantum; a line at a negative frequency
corresponds to stimulated emission into the probe.

The set of lines is a property of the physical Floquet states. Choosing another replica for
the final state, $(\qe_\alpha+\ell\Omega,u_{\alpha,m-\ell})$, relabels $X^{(q)}\to
X^{(q+\ell)}$ but leaves every line frequency and strength unchanged. For a Hermitian probe,
$\overline{X^{(q)}_{\alpha\beta}}=X^{(-q)}_{\beta\alpha}$. Without a drive only $q=0$
survives, and $X^{(0)}$ reduces to the static matrix element.

\paragraph{Sum rule.}
Parseval's theorem and the completeness of the Floquet states give
\begin{equation}\label{eq:sum-rule}
  \sum_\alpha\sum_q\abs{X^{(q)}_{\alpha\beta}}^2=\frac1T\int_0^T\dd t\,
  \matel{\psi_\beta(t)}{X(t)\dg X(t)}{\psi_\beta(t)} ,
\end{equation}
the time-averaged weight of the state $X(t)\ket{\psi_\beta}$. The code reports both sides,
\code{captured} and \code{total}. Their ratio, \code{completeness}, is the fraction of the
spectral weight accounted for by the computed final states. A complete set
(\code{selection="all"}) in the full Fock space of the destination sector saturates the sum rule.

\section{Probes}

Probes are specified by name:
\begin{itemize}
\item \code{current}, the phase derivative $\partial H/\partial\phi$, for a flux or phase probe;
\item \code{charge}, the dot occupation $n_d$, for a gate probe;
\item \code{spin_x}, \code{spin_y}, \code{spin_z} and \code{spin_plus}, \code{spin_minus}, for
spin resonance;
\item \code{d_up}, \code{d_dn} and their adjoints \code{d_up_dagger}, \code{d_dn_dagger},
for tunnelling spectroscopy.
\end{itemize}
Any other named model observable may also be used. Two physical distinctions are important.
\begin{itemize}
\item \emph{Time-dependent probes.} A phase drive changes the current operator itself, so the
probe has harmonics. \pkg{} samples the observables together with the Hamiltonian and uses
these harmonics. A probe given as an explicit operator is treated as time independent.
\item \emph{Tunnelling probes} change the fermion parity and connect Floquet states of
different sectors under the same drive. With $S_X(\omega)$ the line spectrum of the probe
$X$ (defined below), the electron-addition spectrum is $A^+(\omega)=S_{d\dg}(\omega)$ and
the electron-removal spectrum is $A^-(\omega)=S_d(-\omega)$, as for the static spectral
function.
\end{itemize}

\section{Initial states, degeneracies and line shapes}

The observed lines depend on the populations of the Floquet states. Determining these
populations requires a model of relaxation through the environment, which is planned.
At present, spectra are computed for explicitly specified initial states and weights. The default
is the dressed continuation of the static ground state (\code{dressed_state(0)}).

Within a degenerate multiplet the Floquet states can be unitarily rotated. Before constructing
the line spectrum, \pkg{} therefore distributes each initial weight uniformly over its degenerate group and
sums the strengths of degenerate final states, which makes the lines independent of the
basis chosen inside a multiplet. Terms between states of one group at $q=0$ are expectation
values and are omitted from the line spectrum. Terms of one group with $q\ne0$ are \emph{elastic}
lines at $q\Omega$: the probe exchanges energy with the drive without changing the state.
They are kept.

Each line $l$, at frequency $\omega_l=\omega^{(q)}_{\alpha\beta}$, has the strength
$s_l=p_\beta\abs{X^{(q)}_{\alpha\beta}}^2$, where $p_\beta$ is the weight of the initial state.
The lines are broadened with Lorentzian or Gaussian shapes $L$ of half width at half maximum
$w$, which gives the spectrum $S(\omega)=\sum_ls_l\,L(\omega-\omega_l)$
(\code{kind="signed"}, the default). \code{kind="rate"} returns $S(\omega)+S(-\omega)$,
which is proportional to the rate \eqref{eq:floquet-golden} at the probe frequency
$\omega>0$. \code{kind="net"} returns $S(\omega)-S(-\omega)$, absorption minus stimulated
emission, in which the elastic lines cancel. Lines into the discretized continuum need a
width larger than the level spacing; a proper treatment of the continuum is planned.

\begin{incode}
%! omega = np.linspace(0., 1.5, 1501)
\begin{lstlisting}[frame=none,backgroundcolor={},aboveskip=0pt,belowskip=0pt]
options = dict(floquet_cutoff=6, selection="all", method="dense")
states = solve_floquet(periodic, 2, Sector(0, 0), options)
table = states.matrix_elements("charge")              # FloquetTransitions
lines = table.lines(initial=states.dressed_state(0))   # SpectralLines
rate = lines.spectrum(omega, width=.002, kind="rate")
odd = solve_floquet(periodic, 2, Sector(1, 1), dict(floquet_cutoff=6, states=1))
removal = states.matrix_elements("d_up", odd)         # odd state -> even states
\end{lstlisting}
For the junction of \cref{ch:practical}, the complete even set needs a dense Sambe
matrix of dimension $13\times546$, which takes a few minutes of computing time;
\code{basis="eigen"} reduces the cost. The two sets of states must belong to the same drive
and must treat the c-number terms of \cref{sec:gauge} in the same way; their sectors, QP
cutoffs, bases and methods may differ. Static eigenstates also provide
\code{matrix_elements}.
\par\smallskip\usermanual{docs/floquet.md}.
\end{incode}

\section{Computing the harmonics}

Equation \eqref{eq:floquet-element} is a Fourier coefficient of the mode matrix element,
\begin{equation}
  X^{(q)}_{\alpha\beta}=\frac1T\int_0^T\dd t\,\ee^{-\ii q\Omega t}\,
  \matel{u_\alpha(t)}{X(t)}{u_\beta(t)} .
\end{equation}
\pkg{} first forms the harmonics $\ket{v_{\beta,m}}=\sum_kX_k\ket{u_{\beta,m-k}}$ of
$X(t)\ket{u_\beta(t)}$ in the Fock basis of the destination sector. Within one basis this is
a sparse matrix--vector product; between different sectors or cutoffs, the operator acts
between the two QP-truncated spaces. The norm $\sum_m\norm{v_{\beta,m}}^2$ is the
right-hand side of the sum rule \eqref{eq:sum-rule}. The overlap
$\braket{u_\alpha(t)}{v_\beta(t)}$ is a trigonometric polynomial with a finite number $N$ of
harmonics, so $N$ equally spaced samples and a discrete Fourier transform give every
$X^{(q)}$ exactly, without aliasing. The memory is dominated by the vectors $v_{\beta,m}$;
the argument \code{max_memory_gib} limits it, and \code{ket_states} restricts the initial
states.

\section{Mollow triplet and Autler--Townes doublet}\label{sec:mollow}

A resonant pump dresses two levels $a$ and $b$ into the hybrids $\ket\pm$, separated by
the Rabi splitting $\delta$ (\cref{ch:ex-tls}). The probe response depends on whether it acts
on the pumped transition itself or on a transition to a third level.
\begin{itemize}
\item \emph{Probe of the pumped transition, from a dressed state.} The probe connects
$\ket\pm$ with their own replicas and with each other. A system in the lower dressed state
shows lines at three probe frequencies near the pump: absorption at $\Omega+\delta$,
stimulated emission at $\Omega-\delta$, and two elastic lines at the line frequencies
$\pm\Omega$; in the upper dressed state, absorption and emission exchange places. At
resonance each of these four lines carries a quarter of the static strength
$\abs{X_{ab}}^2$, so in the rate the central line is twice as strong as the sidebands. This
is the probe counterpart of the Mollow triplet of resonance
fluorescence~\cite{mollow1969power,mollow1972stimulated,cohentannoudji1998atom}.
The net absorption is negative at $\Omega-\delta$: a dressed two-level system amplifies a
probe at this frequency~\cite{wu1977observation}.
\item \emph{Probe of a transition from a spectator level $c$.} The single line $c\to a$
splits into two lines at $\pm\delta/2$ around its undriven position, each with half the
strength: the Autler--Townes doublet~\cite{autler1955stark}.
\end{itemize}
Both were observed together in a pumped superconducting qubit~\cite{baur2009measurement}.
They appear in \cref{ch:ex-twotone} for the pair transition of the junction, and in
\cref{sec:esr-probe} for the spin doublet.
